% Figure: compartment diagram (stocks and flows) for the SIR model, with
% the resulting S, I, R trajectories alongside.
\begin{figure}[htbp]
\centering
\begin{tikzpicture}[font=\small,
  stock/.style={draw=engblack!75, fill=engblue!12, minimum width=1.2cm,
    minimum height=0.9cm, font=\small\bfseries},
  flow/.style={->, thick, draw=engblack!70}]
% ===== Top: compartment diagram =====
\begin{scope}[yshift=0cm]
\node[stock] (S) at (0,0) {$S$};
\node[stock] (I) at (3,0) {$I$};
\node[stock] (R) at (6,0) {$R$};
\draw[flow] (S) -- node[above, font=\scriptsize] {$\beta SI/N$} (I);
\draw[flow] (I) -- node[above, font=\scriptsize] {$\gamma I$} (R);
\node[font=\scriptsize, anchor=north] at (S.south) {susceptible};
\node[font=\scriptsize, anchor=north] at (I.south) {infectious};
\node[font=\scriptsize, anchor=north] at (R.south) {removed};
\node[font=\small, anchor=south west] at (-1.0, 0.9) {\textbf{(a) stocks and flows}};
\end{scope}
% ===== Bottom: trajectories =====
\begin{scope}[yshift=-4.2cm]
\draw[->, thick, draw=engblack!70] (0,0) -- (8.5,0)
  node[right, font=\scriptsize] {time $t$ (days)};
\draw[->, thick, draw=engblack!70] (0,0) -- (0,3.6)
  node[above, font=\scriptsize] {fraction of $N$};
\foreach \yt/\yv in {0/0, 1.5/0.5, 3.0/1.0} {
  \draw[draw=engblack!70] (0,\yt) -- (-0.07,\yt);
  \node[font=\scriptsize, anchor=east] at (-0.09,\yt) {\yv};
}
% S(t): high then falls
\draw[very thick, draw=engblue!80, smooth, samples=80, domain=0:8]
  plot (\x, { 3.0/(1 + exp(1.4*(\x-3.4))) * 0.95 + 0.05*3.0 });
% I(t): bump
\draw[very thick, draw=engred!85, smooth, samples=80, domain=0:8]
  plot (\x, { 3.0*0.42*exp(-((\x-3.6)*(\x-3.6))/1.6) });
% R(t): rises to plateau
\draw[very thick, draw=engblack!65, smooth, samples=80, domain=0:8]
  plot (\x, { 3.0*0.9/(1 + exp(-1.4*(\x-3.8))) });
\draw[very thick, draw=engblue!80] (5.6,3.2) -- (5.95,3.2);
\node[anchor=west, font=\scriptsize] at (6.0, 3.2) {$S$};
\draw[very thick, draw=engred!85] (5.6,2.8) -- (5.95,2.8);
\node[anchor=west, font=\scriptsize] at (6.0, 2.8) {$I$};
\draw[very thick, draw=engblack!65] (5.6,2.4) -- (5.95,2.4);
\node[anchor=west, font=\scriptsize] at (6.0, 2.4) {$R$};
\node[font=\small, anchor=south west] at (0.2, 3.4) {\textbf{(b) trajectories}};
\draw[dashed, draw=engblack!45] (3.6,0) -- (3.6, 1.5)
  node[above, font=\scriptsize, anchor=south] {peak $I$};
\end{scope}
\end{tikzpicture}
\caption[The SIR compartment model: stocks, flows, and
trajectories]{(a) The susceptible-infectious-removed (SIR) model drawn as
stocks and flows. Each compartment is a stock, an accumulation measured
in head-count; each arrow is a flow, a rate measured in head-count per
unit time. The infection flow $\beta S I / N$ is a reinforcing coupling
(more infectious cases drive more new infections); the removal flow
$\gamma I$ is a draining loop. (b) The trajectories that the coupled rate
equations produce: $S$ falls monotonically, $I$ rises to a peak then
falls, $R$ accumulates toward the final epidemic size. The peak in $I$ is
the engineering quantity a hospital planner needs; its height and timing
are set by the basic reproduction number $R_0 = \beta/\gamma$, the single
dimensionless group that governs the model. Kermack and McKendrick
derived the structure in 1927~\cite{paper:v6c08-kermack-mckendrick-1927}.}
\label{fig:vol02:ch18:compartment-sir}
\end{figure}
